Language models often abandon a correct answer when a user pushes back without giving a reason, and this behavior is usually studied as sycophancy to be suppressed. The same pushback, however, also leads models to revisit and sometimes repair a wrong answer. We measure both effects on multiple-choice questions from MMLU and ARC-Challenge with four open models from two families (Qwen3 1.7B, 4B and 8B, and OLMo-2 7B), reading the final answer exactly from option logits after a free-form reply. The reply budget changes the verdict. When Qwen3-4B may answer the challenge in at most 96 tokens, pushback leaves accuracy unchanged or slightly lower; with 1,024 tokens, it corrects 43\% of its wrong answers and pushback raises accuracy by 5.7 points. Whether pushback helps overall follows a simple break-even rule in the cave and correction rates. Inside the model, caving can be partly predicted before the challenge: hidden states add about 0.04--0.06 AUROC over turn-1 confidence in three of our four probing runs. A probe trained only on self-correction cases predicts caving beyond confidence when replies are truncated, but not when the model can finish its reasoning. Activation steering never reduced caving selectively: where it worked, it suppressed self-correction as well. We conclude that sycophancy evaluations should allow full replies and report self-correction alongside caving.
